% Moderne Dokumentklassen-Definition für LuaLaTeX
% KOMA-Script Report optimiert für LuaLaTeX

\documentclass[
  12pt,                      % Größere Schrift (besser lesbar, DIN 5008)
  a4paper,                   % Papierformat A4
  twoside,                   % Doppelseitiger Druck
  openright,                 % Neue Kapitel auf rechter Seite
  bibliography=totoc,        % Literaturverzeichnis im Inhaltsverzeichnis
  listof=totoc,              % Verzeichnisse im Inhaltsverzeichnis
  index=totoc,               % Index im Inhaltsverzeichnis (falls verwendet)
  tablecaptionabove,         % Tabellenüberschrift oben
  headinclude,               % Kopfzeile im Satzspiegel
  footinclude=false,         % Fußzeile außerhalb (Standard)
  numbers=noenddot,          % Keine Punkte nach Nummern (modern)
  captions=tableheading,     % Optimiert für Tabellen-Captions
  titlepage=true,            % Separate Titelseite
  abstract=true,             % Abstract-Umgebung verfügbar
  version=last,              % Neueste KOMA-Version (MUSS hier stehen!)
  final                      % final Final-Version (keine Draft-Markierungen)
]{scrreprt}

% --- KOMA-Script Optionen (ohne version) ---
\KOMAoptions{
  parskip=half,              % Halber Absatzabstand
  chapterprefix=true,        % "Kapitel X" als Präfix
  headsepline=true,          % Linie unter Kopfzeile
  footsepline=false,         % Keine Linie über Fußzeile
  % Moderne Typografie-Optionen
  DIV=12,                    % Satzspiegel-Berechnung (guter Kompromiss)
  BCOR=0mm,                  % Bindekorrektur (anpassen falls Bindung)
  % Überschriften-Stil
  headings=normal,           % Standard-Überschriften
  % Fußnoten
  footnotes=multiple         % Mehrere Fußnoten pro Seite möglich
}

% --- Schriftgrößen für verschiedene Elemente ---
\addtokomafont{disposition}{\rmfamily\bfseries}  % Überschriften serifenlos + fett
\addtokomafont{chapter}{\large}                  % Kapitel etwas größer
\addtokomafont{section}{\normalsize}             % Abschnitte normal
\addtokomafont{subsection}{\small}               % Unterabschnitte kleiner
\addtokomafont{caption}{\small\itshape}          % Bildunterschriften klein + kursiv
\addtokomafont{captionlabel}{\bfseries}          % "Abbildung X:" fett

% --- Seitenränder (moderne Werte) ---
% Alternative zu DIV-Berechnung - falls manuelle Kontrolle gewünscht
% \usepackage[
%   top=2.5cm,
%   bottom=2.5cm,
%   left=3cm,
%   right=2.5cm,
%   headheight=1.2cm,
%   headsep=0.5cm,
%   footskip=1cm
% ]{geometry}

% --- Kapitel-Stil anpassen (optional) ---
\renewcommand*{\chapterformat}{%
  \chapappifchapterprefix{\nobreakspace}\thechapter\autodot\enskip
}

% --- Moderne Nummerierungstiefe ---
\setcounter{secnumdepth}{3}    % Nummerierung bis Unterabschnitt
\setcounter{tocdepth}{3}       % Inhaltsverzeichnis bis Unterabschnitt

% --- PDF-Kompatibilität ---
% Automatische Erkennung der LaTeX-Engine
\usepackage{iftex}
\ifLuaTeX
  % LuaTeX-spezifische Einstellungen
  \RequirePackage{luatex85}    % Rückwärts-Kompatibilität
\else
  \PackageError{documentclass}{Diese Vorlage ist für LuaLaTeX optimiert!}{%
    Bitte verwenden Sie LuaLaTeX als Compiler.
  }
\fi

% --- Debug-Information ---
\typeout{=== Dokumentklasse geladen ===}
\typeout{Engine: LuaTeX}
\typeout{Klasse: scrreprt (KOMA-Script)}
\typeout{Schriftgröße: 12pt}
\typeout{Papier: A4, zweiseitig}
\typeout{==============================}